\section{Basis Completeness: Identity Is Provably Inert}\label{sec:completeness}

The trunk carries the parameters; this section shows that, at the paper's default rank, the basis
carries nothing the learned projection cannot undo. When the forecast head is linear --- the
default configuration, without the $\dagger$ extension of Table~\ref{tab:legend} --- no
nonlinearity sits between the learned projection and the frozen basis. At $\df = \Ho$, the
non-tiered default throughout, the forecast basis $\Bf$ is square with orthonormal rows: a
complete orthonormal basis of the horizon space. Two complete orthonormal bases differ only by a
rotation, and the freely learned projection absorbs that rotation exactly.

\smallskip
\noindent\textbf{Proposition 1 (rotation absorption).}\quad\emph{Let the forecast head be linear
and let $\Bf$ be complete ($\df = \Ho$). For any complete orthonormal basis $\Bf'$ and any
learned projection $\mathbf{W}_{f}$, there exists a projection $\mathbf{W}_{f}'$ giving the
identical forecast map under $\Bf'$.}

\smallskip
\noindent\emph{Proof sketch.} $\Bf' = \mathbf{Q}\,\Bf$ for some rotation $\mathbf{Q}$, and
$\mathbf{W}_{f}\Bf^{\top} = (\mathbf{W}_{f}\mathbf{Q}^{\top})(\mathbf{Q}\Bf^{\top})$ for any
rotation $\mathbf{Q}$; set $\mathbf{W}_{f}' = \mathbf{W}_{f}\mathbf{Q}^{\top}$.\qed

\smallskip
We verified the absorption numerically at a residual of $5.6 \times 10^{-15}$. On a complete
forecast path, basis identity cannot change what a block can synthesize. This single observation
accounts for three results the submission reported without explanation: the null across wavelet
families, the null across filter lengths within a family, and the interchangeability of blocks at
equal rank.

\paragraph{The measurement.} The conclusion must also hold in practice, at matched capacity. We
swapped the forecast basis among a compactly supported wavelet (Daubechies-3), a discrete cosine
transform, and two random orthonormal draws, at identical parameter counts, across the six M4
periods at ten seeds: 36 paired comparisons, none significant (minimum $p = 0.084$). Read against
the identity disclosure below, this family null is restated over genuinely distinct families
only, as Table~\ref{tab:identity} makes precise.

\paragraph{Scope: expressiveness, not trainability.} The rotation argument fixes what the
forecast head can represent, not whether training finds it. At complete rank every basis is
exactly expressive, yet some complete bases fail to train on some data; the two axes come apart,
and Section~\ref{sec:trainability} measures the second.

\paragraph{Scope: a trend-branch-present null.} Every arm of the swap grid carried a paired
polynomial trend branch, which repairs the one training failure basis identity can cause
(Section~\ref{sec:trainability}). The measured null is therefore a statement about
trend-branch-present configurations, not about frozen-basis stacks in general.

\paragraph{Scope: the forecast path only.} Every shipped backcast basis is truncated to $\Ho$ or
$2\Ho$ rows against a backcast of length between $3\Ho$ and $5\Ho$, and the tiered configurations
truncate the forecast basis as well. Truncated bases drawn from different families span different
subspaces that no rotation relates; on those paths, family equivalence rests on the measurement
rather than on Proposition~1.

\paragraph{Scope: non-tiered configurations only.} Tiered cells discard leading basis rows and so
fall outside the premise of completeness. They belong to the truncated regime of
Section~\ref{sec:truncation}, and nothing in this section speaks to them.
